\documentclass[a4paper,12pt]{article}
\usepackage[margin=2cm]{geometry}
\usepackage[utf8]{inputenc}
\usepackage{hyperref}

\begin{document}

\noindent
{\large \textbf{Štěpán Kovačič}} \hfill {Uživatelská dokumentace k závěrečnému projektu z předmětu NMIN112}\\[-0.2em]

\vspace{0.4em}
\hrule
\vspace{0.6em}

\section*{Úvod}
Tenhle dokument popisuje, jak si aplikaci stáhnout, spustit a jak se s ní
pak pracuje z pohledu běžného uživatele.

\section{Stažení a spuštění}

Nejdřív je potřeba naklonovat repozitář:

\begin{verbatim}
git clone https://github.com/stepakovacic/ZaverecnyProjektMFFProg.git
cd ZaverecnyProjektMFFProg
\end{verbatim}

Dál je potřeba nainstalovat závislosti ze souboru \texttt{requirements.txt}:

\begin{verbatim}
pip install -r setup/requirements.txt
\end{verbatim}

Pro export faktur do PDF musí mít počítač taky nainstalovaný
\texttt{pdflatex}. Bez toho export
sice zbytek aplikace neovlivní, ale tlačítko "pdf export" nebude fungovat.

Nakonec se aplikace spustí příkazem (z adresáře \texttt{app}):

\begin{verbatim}
cd app
uvicorn main:App --reload
\end{verbatim}

Aplikace pak poběží na adrese \url{http://127.0.0.1:8000}, kterou stačí
otevřít v prohlížeči.

\section{Registrace a přihlášení}
Na hlavní stránce je odkaz na registraci. Stačí vyplnit uživatelské jméno,
email a heslo - zbytek (jméno, IČO, telefon, číslo účtu) je nepovinné a dá se
doplnit později přes stránku "Profil". Po registraci se člověk přesměruje na
přihlašovací stránku, kde zadá stejné jméno a heslo, který si při registraci
zvolil.

Po úspěšném přihlášení zůstává uživatel přihlášený i po zavření a znovu
otevření prohlížeče, dokud se sám neodhlásí přes odkaz v navigaci.

\section{Vytvoření faktury}
V navigaci je odkaz "Přehled faktur", kde jde vidět seznam všech faktur
patřících přihlášenému uživateli. Nová faktura se vytváří přes formulář, kde
se vyplní údaje o odběrateli (jméno, adresa, IČO) a datumy vystavení a
splatnosti. Číslo faktury se generuje automaticky (formát rok + pořadové
číslo, např. 2026001), takže ho není potřeba zadávat ručně.

\section{Položky faktury}
Na stránce konkrétní faktury jde přidávat jednotlivé položky - název,
množství a cena za kus. Celková cena položky i celkový součet faktury se
počítá automaticky. Položku jde i smazat, pokud se přidala omylem.

\section{Štítky}
Uživatel si může vytvořit vlastní štítky (název + barva) přes stránku
"Tagy". Vytvořený štítek jde pak přiřadit k libovolné faktuře - jedna
faktura může mít víc štítků a stejný štítek jde použít na víc faktur
zaroveň. Smazání faktury štítek nesmaže, protože štítek patří uživateli, ne
konkrétní faktuře.

\section{Export do PDF}
Na stránce faktury je tlačítko "pdf export", který vygeneruje fakturu ve
formátu PDF a nabídne ji ke stažení. Pokud se faktura od posledního stažení
změnila (přidala/smazala se položka), PDF se při dalším stažení vygeneruje
znovu s aktuálními daty.

\section{Odhlášení}
Odkaz "Odhlásit se" v navigaci ukončí přihlášení a přesměruje zpět na
přihlašovací stránku.

\end{document}